Our approach consists of three sequential stages: activation extraction from a pre-trained ViT backbone, training a Sparse Autoencoder on those activations, and evaluating the learned features using CLIP as an external cross-modal labelling tool. Each stage is described below.

\subsection{Backbone Model and Dataset}

We use a pre-trained ViT-B/16 \cite{Dosovitskiy2021} as the model under analysis. The architecture divides each input image into a sequence of non-overlapping $16 \times 16$ pixel patches, projects them into a $d=768$-dimensional embedding space, and processes them through 12 transformer blocks. A learnable \texttt{[CLS]} token is prepended to the patch sequence; after the final layer, its representation is passed to a linear classification head. The model is used in its standard supervised form, trained on ImageNet-1k.

For activation extraction we use a subset of the ImageNet-1k validation set, covering all 1000 classes. Using the full validation set rather than training data avoids any confound between the features discovered by the SAE and the specific images used to train the ViT backbone.

\subsection{Activation Extraction}

Forward hooks are registered on the outputs of the MLP sub-blocks at layers $\{4, 8, 12\}$, corresponding to the early, middle, and final thirds of the network. For each image, the MLP output at each monitored layer produces a matrix of shape $(N_\text{patches} + 1) \times d$, where $N_\text{patches} = 196$ for a $224 \times 224$ input. We retain only the patch token activations (excluding the \texttt{[CLS]} token at this stage), yielding per-image activation matrices of shape $196 \times 768$ per layer.

All activations are collected, centred by subtracting the dataset-level mean, and stored on disk prior to SAE training. This pre-processing step is standard practice and has been shown to improve the stability and interpretability of the learned dictionary \cite{Cunningham2024}.

\subsection{Sparse Autoencoder}

The SAE consists of a single linear encoder followed by a ReLU non-linearity, and a linear decoder constrained to have unit-norm columns. Given an activation vector $\mathbf{x} \in \mathbb{R}^d$, the encoder produces a sparse latent code:
\begin{equation}
    \mathbf{z} = \text{ReLU}\!\left(\mathbf{W}_\text{enc}\,\mathbf{x} + \mathbf{b}_\text{enc}\right), \quad \mathbf{z} \in \mathbb{R}^m,
\end{equation}
where $m \gg d$ is the dictionary size (we use $m = 4d = 3072$). The decoder reconstructs the original activation as $\hat{\mathbf{x}} = \mathbf{W}_\text{dec}\,\mathbf{z}$. The training objective combines a reconstruction loss with an $L_1$ sparsity penalty:
\begin{equation}
    \mathcal{L} = \underbrace{\|\mathbf{x} - \hat{\mathbf{x}}\|_2^2}_{\text{reconstruction}} + \lambda \underbrace{\|\mathbf{z}\|_1}_{\text{sparsity}},
\end{equation}
where $\lambda$ controls the trade-off between reconstruction fidelity and feature sparsity. A separate SAE is trained for each monitored layer, allowing us to compare the character of features across depth.

\subsection{Cross-Modal Feature Evaluation via CLIP}

Once training is complete, each of the $m$ SAE features is a direction $\mathbf{d}_k = \mathbf{W}_\text{dec}[:, k]$ in the original activation space. To assign a human-readable label to each feature, we identify the top-$K$ images (we use $K=16$) that most strongly activate feature $k$, and use CLIP \cite{Radford2021} to match those images against a textual vocabulary.

Concretely, we compute the CLIP image embeddings of the top-$K$ images and the CLIP text embeddings of a vocabulary of candidate labels (ImageNet class names augmented with a set of texture and colour descriptors). The label assigned to feature $k$ is the text entry with the highest mean cosine similarity across the $K$ images. This procedure is performed independently for each feature and each layer, producing a labelled dictionary of visual concepts at three levels of network depth.

A feature is considered \textit{monosemantic} if the same label (or a close semantic neighbour) accounts for the top-activating images. Features whose top-activating images span unrelated categories are flagged as polysemantic and excluded from the interpretability analysis, providing a natural quality filter on the discovered concepts.

% TODO: fill in lambda, lr, batch size, epochs, hardware.
